\section{Edition and integration choices for the lending query}
\label{exec14:commercial-contracts}

An actionable comparison must connect the banking request to the particular
operation that an edition exposes. The population effect of a repayment offer,
the effect transmitted through payment behavior, and an alternative outcome
for an already observed borrower should therefore have different request
contracts. More detailed Bayes Server documentation supplies a useful test of
this distinction. Its Advanced Causal Query guide describes back-door and
front-door criteria, subsequent adjustment calculations, multiple treatments
and joint outcomes. Support for hidden confounders and direct effects varies
with the selected algorithm. Those are documented capabilities with which the
proposal must compete; the list is not evidence of a complete general
identification algorithm \citep{exec14_ci_causal_query}.

The version evidence has several layers. The guides describe the features as
available from version~10. The acquired release history associates advanced
causal queries, back-door and front-door algorithms, causal effects analysis
and the abduction helper with version~10.2 beta, dated 19 June 2022. Its highest
listed release is 11.4.0.0, dated 17 April 2026. These observations support a
version-family comparison; they do not identify the bytes of an installed
assembly or establish equivalence across its interfaces
\citep{exec14_ci_release_history}. The copies were acquired on 14 October
2026. As elsewhere in this survey, dated release entries and copyright notices
do not authenticate a mutable page's exact contents on the 7 October cutoff.

\subsection{A narrow identification adapter has a concrete contract}

The documented \texttt{FrontDoorCriterion} class takes a supplied network. Its
\texttt{Identify} methods receive interventions, outcomes, other evidence
and options, and return criterion-specific identification information without
performing estimation. Associated methods, \texttt{IdentifyXZ} and
\texttt{IdentifyZY}, address the adjustment subproblems between treatment,
front-door nodes and outcomes. The documentation explicitly envisages using
the identified node sets outside Bayes Server. It also distinguishes a Boolean
validity check from validation that raises an exception
\citep{exec14_ci_api_frontdoor}. This is a much more useful integration lead
than a claim that a product ``supports causality.''

For the repayment example, a proposed adapter would receive the offer variable
$A$, a mediator $M$, repayment $Y$, an observed-variable map, a supplied causal
network and the desired intervention. It would return the criterion, the
identified node sets, the input model identity and the query to which they
apply. An external estimator would then require the observed distributions
and support conditions needed to evaluate the functional. The exact
front-door calculation developed earlier illustrates that separation:
returning a mediator name is not the same as fitting the required conditional
distributions, and fitting them is not proof that the graph describes the bank.

A useful minimal adapter could therefore export a reviewed identification
result before undertaking a much larger attempt to embed an entire modeling
environment. The export must preserve node identity, intervention versus
ordinary evidence, treatment coding, relevant clock and units. Its response
should distinguish a valid criterion result, invalid input, unsupported
options and failure to find a set. In particular, failure of this one
criterion must not become a universal nonidentification certificate. That
conclusion would need a complete algorithm for the declared graph/query class
or an independent argument. The selected documentation establishes neither
general ID/IDC completeness nor unrestricted counterfactual identification.

The reference revision examined in this comparison exposed supplied-DAG
back-door checks, supplied-moment Gaussian residual inversion and explicit
sensitivity and clock operations. It did not expose this front-door adapter,
a general abduction engine or a fitted mediation service. Later fitted-summary
inference in Chapter~\ref{gap29:public} is separate evidence and does not
qualify those missing commercial interfaces. The proposed
scientific contract can organize those future capabilities, but it does not
erase the present functional gap. An investment case should demonstrate the
value of preserving scientific meaning across an existing estimator and a
specific identification service before proposing another broad implementation. \citep{exec14_ci_comparison}

\subsection{Effect labels need mathematical definitions}

The Advanced Causal Query guide describes total and direct effects using paths.
One illustrative total-effect list includes a path of the form
$A\leftarrow C\rightarrow Y$. That is a back-door association path, not a
directed causal route from $A$ to $Y$. This wording should not be adopted as
the definition of total effect, and it does not by itself demonstrate an
implementation error \citep{exec14_ci_causal_query}. A wrapper must bind the
requested quantity to an explicit mathematical expression and then determine
whether the selected vendor operation computes it.

For example, the total contrast is
$\mathbb E[Y(1)-Y(0)]$. A controlled direct contrast at mediator level $m$ is
$\mathbb E[Y(1,m)-Y(0,m)]$, while a natural direct contrast uses the
unit-specific mediator $M(0)$ inside a nested potential outcome. A menu label
about direct links does not settle which of these interventions and
cross-world premises is intended. Stata's specified mediation workflow
provides a useful contrast because it states its natural-effect targets and
confounding restrictions \citep{exec13_ci_stata_mediate}. The proposal should
compare matching estimands, not merely similarly named report columns.

\section{Abduction must preserve uncertainty about the borrower}

The Bayes Server abduction dialog describes querying characteristic variables
given factual evidence, removing evidence for the individual and setting
evidence on those characteristics. Its .NET \texttt{Abduction.Update} signature
separates current evidence, the variables supplying abduction evidence, the
characteristic variables and additional options. This is explicit support for
an abduction step. The selected pages do not fully specify how an uncertain
multivariate posterior is represented or propagated through every supported
case \citep{exec14_ci_abduction,exec14_ci_api_abduction}. That remaining
question matters even when the conditional mean of each characteristic can
be computed accurately.

Let $U$ represent latent repayment capacity in normalized units. For an
illustrative required payment burden $a$, define
\[
 U\sim\operatorname{Uniform}\{0,1,3\},\qquad
 Y(a)=\mathbf 1\{U\geq a\}.
\]
Assignment of the burden is independent of $U$ in this stipulated model.
For a borrower observed at $A=1$ with repayment $Y=1$, Bayes' rule eliminates
$U=0$ and retains $U=1$ and $U=3$ with equal posterior mass. Under the
alternative burden $a'=2$, the individual-conditional query is
\[
 \Pr\{Y(2)=1\mid A=1,Y=1\}
 =\sum_{u\in\{1,3\}}\mathbf 1\{u\geq2\}
   \Pr(U=u\mid A=1,Y=1)=\frac12.
\]
The population intervention probability is instead $1/3$. Replacing the
posterior of $U$ by its mean, $\mathbb E[U\mid A=1,Y=1]=2$, gives the
indicator $\mathbf 1\{2\geq2\}=1$. All three calculations concern the
same structural equation, but they answer different questions or use a lossy
approximation. The finite enumeration reproduces the fractions exactly; it
is a project construction, not a test of Bayes Server
\citep{exec14_ci_exact}.

More generally, a model-specified individual counterfactual averages the
alternative structural response over the posterior of the same exogenous
state that generated the factual observations. Applying a nonlinear response
to a posterior mean need not equal that average. For several exogenous
variables, preserving only separate marginal distributions can also discard
dependence needed by the response. A capability contract should state whether
the operation returns a uniquely recovered characteristic value, a posterior
distribution, a point approximation or a set of compatible states. It must
also identify the extra structural assumptions that connect alternative
outcomes for the same borrower.

To make the financial consequence explicit, stipulate a USD~1,000 payment
upon success and a USD~20 cost under the alternative scenario. The posterior
calculation gives an expected payoff of $1{,}000(1/2)-20=480$ dollars. The
mean substitution gives $1{,}000(1)-20=980$ dollars, an overstatement of
USD~500 in this model. These are alternative-scenario payoffs, not incremental
policy profit: a decision comparison would additionally specify the baseline
payoff and all costs. Neither the amounts nor the latent law are bank
measurements \citep{exec14_ci_exact}.

The immediate public/generated test is therefore precise. Supply a small
structural model whose exogenous posterior has more than one possible state,
compare its exact enumerated counterfactual distribution with the returned
object, and repeat with dependent exogenous variables. Test removal of factual
evidence separately from the intervention. Retain unsupported or approximate
behavior explicitly. The inspected signature is sufficient to motivate this
prototype experiment; it is insufficient to declare the result in advance.

\section{Sweeps, units and uncertainty in the commercial output}

The Effects Analysis guide describes sweeps over intervention values. Discrete
treatments use their states; continuous treatments are discretized conditional
on current evidence. For discrete outcomes the requested state probability is
reported; continuous outcomes can produce a mean and variance. The guide also
normalizes the chart's horizontal axes when treatments are compared
\citep{exec14_ci_effects}. Those choices affect what an exported plot means.

A lending manager should be able to recover the actual payment burden or
interest-rate change corresponding to every displayed point. A normalized
coordinate without the original dose, unit and transformation is not an
actionable pricing instruction. If the grid depends on factual evidence, a
different borrower's curve may use different underlying values. A prospective
adapter must carry the dose grid and conditioning evidence as data, rather
than reconstructing them from an image. The intervention guide further
describes discrete, continuous and dynamic-network interventions, which must
remain distinct from ordinary observational evidence
\citep{exec14_ci_interventions}.

The reported variance of an outcome under a fixed fitted model is also a
different object from uncertainty about its estimated causal mean. If
$\theta$ denotes parameters fitted from data, then
$\operatorname{Var}_{\theta}\{Y(a)\}$ describes outcome variability at
that model, while sampling variation of an estimator of
$\mathbb E_{\theta}\{Y(a)\}$ requires an estimation and sampling law.
Uncertainty about the graph or about an unmeasured causal premise is another
layer. A mean/variance pair from an intervention sweep cannot be relabeled as
a calibrated confidence interval for the offer effect. The proposed result
format should identify these layers and report which are included.

\section{Reuse choices across the compared families}

BayesiaLab's supplied-graph and data-mapping workflow remains a credible
commercial comparator. Its Engine API overview distinguishes Java modeling,
learning and inference configurations from REST inference. The table places
Java model files on customer servers and REST model files on Bayesia servers
in France; production capacity is tied to permitted concurrent instances,
with documented license-validation connectivity. These facts matter to an
integration design, but they do not establish that every graphical causal
function is exposed by every configuration
\citep{exec13_bayesia_api,exec13_bayesia_professional}. The proposed route is
a narrowly specified graph/query export followed by verified interpretation
of the returned result. The missing causal API mapping must be acquired before
committing to an inference-only endpoint for a more elaborate task.

Stata supplies a mature treatment and mediation baseline for specified targets.
Its treatment-effect documentation is more directly relevant to a one-time
offer ATE or ATT than a generic graph-learning comparison
\citep{exec13_ci_stata_causal_intro}. The license-options text distinguishes
an identified user's installations from concurrent network use at a physical
site. This is useful deployment context, but the options page does not settle
the entitlement for a proposed multiuser service
\citep{exec13_ivrd_license}. A first comparison can use an appropriately
authorized analyst workflow and complete output artifacts. Any subsequent
CLI/MCP service needs a separately resolved execution and export contract.

The retained JMP 19.2 documentation supplies the treatment comparison that
the initial acquisition had left unresolved. The JMP Pro Causal Treatment personality describes ATE and
ATET for binary treatment, and potential-outcome means for continuous treatment.
Its assumptions distinguish consistency, positivity and conditional
exchangeability. Its binary estimators use specified propensity and outcome
models, with weighting, regression adjustment and augmentation alternatives
\citep{exec13_ch_jmp_causal}. These are meaningful commercial capabilities
that the local reference runtime does not yet provide.

The documentation's potential-outcome predictions should not be described as
identified missing individual outcomes. Model-based conditional means and
population averages have their own legitimate roles; they do not supply the
cross-world dependence in the borrower example. A proposed file-based
comparison should preserve target weights, observation exclusions, nuisance
specifications and treatment coding before attempting unattended automation.
Only inspected prose supports this comparison; image-only equations were not
treated as inspected technical evidence. Installed build, automation rights
and selected deployment prices remain unresolved.

BayesFusion's SMILE documentation provides a compiled C++ library and named
wrappers, as well as separate development and deployment licensing categories
\citep{exec13_smile_intro,exec13_smile_license}. That is adjacent integration
evidence. Section~\ref{sec:exec24-smile-reuse} develops the supplied-model
decision and temporal interfaces; a matched general intervention/counterfactual
contract remains a separate qualification.
For HUGIN, Netica, SAS, IBM, Mplus, MathWorks, causaLens, EViews, OxMetrics,
LISREL and Analytica, the technical or edition-specific questions retained
earlier in this chapter remain open at different levels. Some have useful
adjacent manuals or original-paper leads; inaccessible or incomplete pages
do not establish absent functionality. Broader ranking should wait for the
matching source rather than borrow a conclusion from another software area. \citep{exec14_ci_comparison}

\section{Published prices and conditions for a justified comparison}

The acquired Bayes Server purchase page explicitly selects Commercial and GBP
in its form. It lists version-11 Professional individual UI and API developer
licenses at GBP~550 each, and their combination at GBP~1,100, before applicable
VAT. It describes perpetual licenses for the major version with minor upgrades
and per-developer API functionality within the purchaser's own product
\citep{exec14_ci_purchase}. These are public page observations on
14 October 2026, not a quotation, a historical-cutoff price or total bank
deployment cost. They should remain in their original currency; conversion
would need an explicit exchange-rate date and cost purpose.

The linked, undated software agreement contains restricted executable-use and
transfer terms and a requirement for prior written consent before disclosing
benchmark results to third parties. It allows that an amendment or addendum
may accompany the software \citep{exec14_ci_eula}. This does not justify a
blanket conclusion that the purchase-page description is false or that a
particular proposed distribution is licensed. It identifies an exact remaining
commercial question: reconcile the selected order and applicable agreement
before promising embedded deployment or public benchmark publication. No
license was accepted, no vendor was contacted and no commercial binary was
executed for this document comparison.

The proposal can proceed with a useful staged decision. First define a narrow
target and compare its documented contract, information requirements and
failure states against the local service. Next build a generated/public-data
prototype that demonstrates the missing semantic preservation and justifies
the cost of a matched experiment. Where a proprietary engine is involved,
resolve the actual execution and reporting terms before claiming a measured
comparison. Finally measure accepted output quality, review time, engineering
effort and operating costs together. A lower license amount, additional menu
entry or successful API response cannot alone decide which route gives the
bank a more defensible answer.
\label{exec14:commercial-contracts-end}
